% Generated draft: requires mathematical and source review.
\documentclass[11pt]{article}
\usepackage[margin=25mm]{geometry}
\usepackage{fontspec}
\setmainfont{Latin Modern Roman}
\usepackage{amsmath,amssymb,mathtools,unicode-math}
\setmathfont{Latin Modern Math}
\usepackage{enumitem,needspace,xurl,hyperref}
\setlength{\parindent}{0pt}
\setlength{\parskip}{5pt}
\setlength{\emergencystretch}{3em}
\newcommand{\N}{\mathbb{N}}
\newcommand{\Z}{\mathbb{Z}}
\newcommand{\Q}{\mathbb{Q}}
\newcommand{\R}{\mathbb{R}}
\newcommand{\C}{\mathbb{C}}
\newcommand{\F}{\mathbb{F}}
\newcommand{\A}{\mathbb{A}}
\newcommand{\PP}{\mathbb{P}}
\newcommand{\E}{\mathbb{E}}
\newcommand{\eps}{\varepsilon}
\newcommand{\dd}{\,\mathrm{d}}
\begin{document}
\section*{UnsolvedMath ID: 9700004}
\textbf{Source problem code:} catalogue entry
\subsection{Definitions and mathematical statement}
A standard finite coagulation--fragmentation model has states describing partitions of a fixed finite mass or set, with transitions that merge blocks or split a block. The rate kernel and splitting rule are model parameters. The source does not specify labeled versus size partitions or a particular family of rates.

For an explicit continuous-time formulation, let \(Q=(q(x,y))\) be the generator of an irreducible chain on a finite state space \(\Omega\), with stationary probability \(\pi\), and assume reversibility,
\[
 \pi(x)q(x,y)=\pi(y)q(y,x).
\]
Let \(P_t=e^{tQ}\), let \(\lambda_1>0\) be the smallest nonzero eigenvalue of \(-Q\), and define
\[
 t_{\mathrm{rel}}=\lambda_1^{-1},\qquad
 t_{\mathrm{mix}}(\varepsilon)=\inf\left\{t\ge0:
 \max_x\frac12\sum_{y\in\Omega}|P_t(x,y)-\pi(y)|\le\varepsilon\right\},
 \quad0<\varepsilon<1.
\]
Obtain useful bounds on these quantities in terms of the intended coagulation--fragmentation parameters. These definitions are a conventional candidate; the source does not choose continuous versus discrete time, irreducibility/aperiodicity conventions, or a mixing threshold.

The actual rate family, reversible measure and parameters are missing. A rate normalization is essential because rescaling all rates rescales the time unit. A meaningful bound must also specify its permitted dependence on mass/particle count, rate variation and splitting structure. No universal polynomial-time law is asserted for arbitrary reversible chains, and generic finite-chain bounds are not confused with a model-specific sharp analysis.

The supplied summary reports special-case literature but no systematic general parameter-dependent bounds verified. Its cited coalescence review and double-coset work are expressly related background rather than solutions. The source's open presentation and later search assessment remain attributions pending primary-literature review.
\subsection{Short English statement}
Bound spectral relaxation and total-variation mixing for reversible finite coagulation–fragmentation chains using their actual rate and size parameters. The chain family, time normalization and target bound are not specified in the excerpt.
\subsection{Sources}
\begin{itemize}
\item Local registry: \url{data/unsolvedmath/problems.json}, permanent UnsolvedMath ID \(9700004\), source code catalogue entry; supplied statement, background, and research summary, updated 2026-08-06.
\item David Aldous, \emph{Open problems}, item coag\_frag.html, maintained index: \url{https://www.stat.berkeley.edu/~aldous/Research/OP/index.html}. Supplied HTML text with linked-primary-PDF text accessed 2026-07-29; the index is reported to have 23 active problem pages.
\item Aldous, ``Deterministic and stochastic models for coalescence (aggregation and coagulation): a review of the mean-field theory for probabilists'', \emph{Bernoulli} (1999), supplied as background rather than a mixing-time result.
\item ``Double Coset Markov Chains'', arXiv (2021), supplied as related work without authors, identifier or a resolution of this question.
\end{itemize}
\end{document}
